%% DARTSim / RS-DRL — Question & Answer Preparation Guide
%% Compiled from Dartsim_questions.md (corrected against implementation)
%% Engine: XeLaTeX  |  Encoding: UTF-8
\documentclass[oneside,openany,a4paper,12pt]{article}

% ─── Math & general packages ─────────────────────────────────────────────────
\usepackage{amsthm,amssymb,amsmath}
\usepackage[top=25mm, bottom=25mm, left=25mm, right=30mm]{geometry}
\usepackage{graphicx}
\usepackage{setspace}
\usepackage{booktabs}
\usepackage{array}
\usepackage{tabularx}
\usepackage{longtable}
\usepackage{enumitem}

% ─── Colors & Boxes ──────────────────────────────────────────────────────────
\usepackage{xcolor}
\usepackage{tcolorbox}
\tcbuselibrary{skins,breakable}

% ─── Hyperlinks ──────────────────────────────────────────────────────────────
\usepackage[pagebackref=false,colorlinks,linkcolor=blue,citecolor=blue]{hyperref}

% ─── Persian / Bidi — same approach as thesis ────────────────────────────────
\usepackage{xepersian}
\settextfont[Path=font/,Extension=.ttf,Scale=1.1,
    BoldFont=XB NiloofarBd,
    ItalicFont=XB NiloofarIt,
    BoldItalicFont=XB NiloofarBdIt
]{XB Niloofar}
\setlatintextfont[Scale=0.9]{Times New Roman}
\setdigitfont[Path=font/,Extension=.ttf,Scale=1.1,
    BoldFont=XB NiloofarBd,
    ItalicFont=XB NiloofarIt,
    BoldItalicFont=XB NiloofarBdIt
]{XB Niloofar}



% ─── Custom Colours ─────────────────────────────────────────────────────────
\definecolor{hintbg}{RGB}{240,248,255}
\definecolor{hintborder}{RGB}{70,130,180}
\definecolor{warnbg}{RGB}{255,248,220}
\definecolor{warnborder}{RGB}{218,165,32}
\definecolor{fixbg}{RGB}{240,255,240}
\definecolor{fixborder}{RGB}{34,139,34}
\definecolor{qbg}{RGB}{245,245,250}
\definecolor{qborder}{RGB}{100,100,180}

% ─── tcolorbox styles ────────────────────────────────────────────────────────
\newtcolorbox{hintbox}[1][]{
    enhanced, breakable,
    colback=hintbg, colframe=hintborder,
    title={\textbf{راهنما}},
    fonttitle=\bfseries,
    #1
}

\newtcolorbox{warnbox}[1][]{
    enhanced, breakable,
    colback=warnbg, colframe=warnborder,
    title={\textbf{نکته مهم / تصحیح}},
    fonttitle=\bfseries,
    #1
}

\newtcolorbox{questionbox}[2][]{
    enhanced, breakable,
    colback=qbg, colframe=qborder,
    title={\textbf{سؤال #2}},
    fonttitle=\bfseries,
    #1
}

% ─── Commands ────────────────────────────────────────────────────────────────
\newcommand{\rsdrl}{\lr{RS-DRL}}
\newcommand{\dqn}{\lr{DQN}}
\newcommand{\dartsim}{\lr{DARTSim}}
\newcommand{\mapek}{\lr{MAPE-K}}
% \newcommand{\rho} is already defined

% ─────────────────────────────────────────────────────────────────────────────
\begin{document}

\begin{titlepage}
\centering
\vspace*{3cm}
{\huge\bfseries راهنمای جامع آمادگی برای دفاع از پایان‌نامه}\\[1em]
{\Large\bfseries شبیه‌ساز \dartsim{} و روش \rsdrl{}}\\[2em]
{\large سؤالات احتمالی استاد با پاسخ‌های مستند}\\[0.5em]
{\normalsize (مقادیر اصلاح‌شده بر اساس پیاده‌سازی واقعی و کد منبع)}\\[3em]
\hrule height 1.5pt
\vspace{1.5em}
{\large علیرضا کاویانی‌فر}\\[0.5em]
{\normalsize \lr{GitHub: alirezakavianifar/RS-DRL\_DartSim}}\\[3em]
\vfill
{\large \today}
\end{titlepage}

\tableofcontents
\clearpage

%% ═══════════════════════════════════════════════════════════════════════════
\section{مفاهیم پایه و انگیزه}
%% ═══════════════════════════════════════════════════════════════════════════

\subsection{فضای مسئله}

\begin{questionbox}{Q1}
چرا \dartsim{} را به عنوان بستر آزمایشی انتخاب کردید؟ چه چالش‌های منحصربه‌فردی ارائه می‌دهد؟
\end{questionbox}

\begin{hintbox}
\begin{itemize}
  \item معامله ارتفاع (\lr{Altitude Trade-off}) به عنوان چالش محوری
  \item $24.4\%$ انتقال‌های تصادفی (عدم قطعیت محیطی)
  \item فضای حالت پیوسته با اقدامات گسسته
  \item الزامات تصمیم‌گیری بلادرنگ ($< 10$ میلی‌ثانیه)
\end{itemize}
\end{hintbox}

\begin{questionbox}{Q2}
«معامله ارتفاع» را با کلمات خود توضیح دهید. چرا چالش «محوری» نامیده می‌شود؟
\end{questionbox}

\begin{hintbox}
تنش بنیادین میان دو هدف متضاد:
\begin{itemize}
  \item \textbf{ارتفاع پایین} $\leftarrow$ شناسایی بهتر اهداف، ولی افزایش ریسک تهدیدات زمینی
  \item \textbf{ارتفاع بالا} $\leftarrow$ کاهش ریسک تهدیدات، ولی کاهش کیفیت شناسایی اهداف
\end{itemize}
\end{hintbox}

\begin{questionbox}{Q3}
تفاوت مسئله \dartsim{} با معیارهای استاندارد \lr{RL} مانند \lr{Atari} یا \lr{MuJoCo} چیست؟
\end{questionbox}

\begin{hintbox}
\begin{itemize}
  \item محدودیت‌های زمان واقعی (بلادرنگ)
  \item بهینه‌سازی چندهدفه (موفقیت مأموریت + بقا + سرعت)
  \item عدم قطعیت بالا ($24.4\%$ انتقال‌های تصادفی)
  \item ماهیت حیاتی‌بودن از نظر ایمنی (انهدام پهپاد)
\end{itemize}
\end{hintbox}

\subsection{یکپارچه‌سازی حلقه \mapek{}}

\begin{questionbox}{Q4}
چرا \mapek{} را با \rsdrl{} یکپارچه کردید؟ این یکپارچه‌سازی چه مشکل خاصی را حل می‌کند؟
\end{questionbox}

\begin{hintbox}
\begin{itemize}
  \item پایش پیوسته و آموزش مجدد بر حسب تقاضا
  \item مقابله با انحراف مفهومی (\lr{Concept Drift}) در محیط‌های پویا
  \item تضمین کیفیت از طریق پایش آستانه‌ها
  \item حلقه بازخورد برای تطبیق‌پذیری
\end{itemize}
\end{hintbox}

\begin{questionbox}{Q5}
اجزای حلقه \mapek{} را در زمینه سیستم شما گام به گام توضیح دهید.
\end{questionbox}

\begin{hintbox}
\begin{description}
  \item[\textbf{\lr{Monitor}}:] معیارهای کلیدی زمان اجرا شامل پاداش، موفقیت مأموریت، اهداف شناسایی‌شده و زمان تصمیم را در قالب شیء \lr{SystemMetrics} جمع‌آوری می‌کند.
  \item[\textbf{\lr{Analyze}}:] معیارها را با آستانه‌های تعریف‌شده در پایگاه دانش مقایسه می‌کند.
  \item[\textbf{\lr{Plan}}:] در صورت نقض آستانه، نقشه وفقی با $5000$ گام، نرخ یادگیری $10^{-4}$ و $\varrho=0.3$ تهیه می‌کند.
  \item[\textbf{\lr{Execute}}:] آموزش مجدد را روی بافر تجربیات آفلاین اجرا می‌کند.
  \item[\textbf{\lr{Knowledge}}:] مدل جاری، تاریخچه آموزش، آستانه‌ها و تاریخچه وفق را ذخیره می‌کند.
\end{description}
\end{hintbox}

\begin{questionbox}{Q6}
چه آستانه‌هایی در حلقه \mapek{} تعریف کردید و چرا آن مقادیر خاص؟
\end{questionbox}

\begin{hintbox}
\textbf{مقادیر تأییدشده از \lr{mape\_k\_architecture.py}:}
\begin{itemize}
  \item حداقل پاداش اپیزودی: $< 0.7$
  \item نرخ موفقیت مأموریت در ۱۰ اپیزود اخیر: $< 0.8$
  \item حداکثر زمان تصمیم‌گیری: \textbf{$1000$ میلی‌ثانیه}
  \item حداقل نسبت اهداف شناسایی‌شده: \textbf{$\geq 0.6$ (۶۰\%)}
  \item قانون «\textbf{۳ نقض متوالی}» قبل از فعال‌شدن آموزش مجدد
  \item حداقل فاصله بین آموزش‌های مجدد: \textbf{۱۰۰۰ گام}
\end{itemize}
\end{hintbox}

%% ═══════════════════════════════════════════════════════════════════════════
\section{روش‌شناسی — عمیق‌سازی}
%% ═══════════════════════════════════════════════════════════════════════════

\subsection{فرمول‌بندی \lr{MDP}}

\begin{questionbox}{Q7}
فضای حالت شما ۱۷ بعد دارد. هر مؤلفه را توضیح دهید و ضرورت هر کدام را بیان کنید.
\end{questionbox}

\begin{hintbox}
\begin{itemize}
  \item \textbf{موقعیت مکانی ($2$ بعد):} مختصات $(X, Y)$ نرمال‌شده
  \item \textbf{جهت حرکت ($2$ بعد):} بردار جهت نرمال‌شده
  \item \textbf{پیکربندی جاری ($5$ بعد):} ارتفاع، آرایش تیمی، \lr{ECM}، \lr{TTC IncAlt}، \lr{TTC DecAlt}
  \item \textbf{حسگر تهدید ($5$ بعد):} بردار باینری ۵ سلول پیشِ‌رو
  \item \textbf{حسگر هدف ($5$ بعد):} بردار باینری ۵ سلول پیشِ‌رو
\end{itemize}
\end{hintbox}

\begin{questionbox}{Q8}
چرا فضای اقدام گسسته (۸ اقدام) را انتخاب کردید نه کنترل پیوسته؟
\end{questionbox}

\begin{hintbox}
\begin{itemize}
  \item سادگی یادگیری
  \item تصمیمات تاکتیکی ذاتاً گسسته هستند
  \item سازگاری با \dqn{}
  \item اما اذعان داشته باشید به محدودیت‌ها: شکاف شبیه‌ساز به واقعیت، تنش مکانیکی
\end{itemize}
\end{hintbox}

\begin{questionbox}{Q9}
تابع پاداش ۴ مؤلفه با وزن‌های خاص دارد. چرا این وزن‌ها ($0.4, 0.3, 0.2, 0.1$)?
\end{questionbox}

\begin{hintbox}
\begin{equation}
R = w_{\text{success}} \cdot r_{\text{success}} + w_{\text{targets}} \cdot r_{\text{targets}} + w_{\text{survival}} \cdot r_{\text{survival}} + w_{\text{efficiency}} \cdot r_{\text{efficiency}}
\end{equation}

\begin{itemize}
  \item $w_{\text{success}} = 0.4$: موفقیت مأموریت مهم‌ترین هدف است
  \item $w_{\text{targets}} = 0.3$: شناسایی اهداف اولویت دوم
  \item $w_{\text{survival}} = 0.2$: بقای تیم؛ در صورت انهدام جریمه سنگین $w_{\text{destruction}} = -0.5$
  \item $w_{\text{efficiency}} = 0.1$: تشویق به تصمیمات سریع‌تر
\end{itemize}
\end{hintbox}

\subsection{نوآوری \rsdrl{}}

\begin{questionbox}{Q10}
\textbf{[مهم‌ترین سؤال]} مکانیزم \rsdrl{} را به تفصیل توضیح دهید. «شکل‌دهی تصادفی پاداش» دقیقاً به چه معناست؟
\end{questionbox}

\begin{hintbox}
\begin{description}
  \item[\textbf{\dqn{} استاندارد}:] از پاداش‌های واقعی $(r)$ می‌آموزد.
  \item[\textbf{\rsdrl{}}:] $K = \lfloor\varrho \cdot N\rfloor$ انتقال از مینی‌بچ که دارای \textbf{پاداش غیرمثبت ($r \leq 0$)} هستند انتخاب شده و پاداش آن‌ها با مقدار خوش‌بینانه $+1.0$ جایگزین می‌شود. \textbf{آستانه پیاده‌سازی: \lr{rewards <= 0.0}} (نه فقط پاداش‌های منفی).
\end{description}

این تغییر باعث می‌شود:
\begin{itemize}
  \item از «انجماد ارزش کیو» جلوگیری شود
  \item به عنوان رگولارایزر پویا عمل کند
  \item خطای \lr{TD Loss} فعال بماند ($0.27$ تا $4.63$ در مقابل $0.006$ در \dqn{})
\end{itemize}

\textbf{نکته:} $\varrho$ \textit{کسر} مینی‌بچ است که بازطراحی می‌شود ($K = \lfloor\varrho \cdot N\rfloor$ انتقال)، نه احتمال هر نمونه.
\end{hintbox}

\begin{questionbox}{Q11}
چرا $\varrho=0.3$ خاص؟ نتایج مطالعه ابلیشن را نشان دهید.
\end{questionbox}

\begin{hintbox}
نتایج جدول مطالعه ابلیشن (سناریوی متوسط):

\begin{center}
\begin{tabular}{ccc}
\toprule
\textbf{$\varrho$} & \textbf{میانگین ارزش کیو} & \textbf{توضیح} \\
\midrule
$0.0$ & $1.948 \pm 0.687$ & معادل \dqn{} پایه \\
$0.1$ & $4.266 \pm 0.826$ & بهبود محدود \\
$\mathbf{0.3}$ & $\mathbf{8.190 \pm 0.605}$ & \textbf{بهینه!} \\
$0.5$ & $12.387 \pm 0.621$ & بیش از حد تهاجمی \\
\bottomrule
\end{tabular}
\end{center}
\end{hintbox}

\begin{questionbox}{Q12}
اگر $\varrho$ خیلی زیاد ($0.5$) یا خیلی کم ($0.0$) باشد چه اتفاقی می‌افتد؟
\end{questionbox}

\begin{hintbox}
\begin{itemize}
  \item $\varrho=0.0$: گیر افتادن در نقاط بهینه محلی، انجماد ارزش کیو، سیاست‌های محافظه‌کارانه
  \item $\varrho=0.5$: خوش‌بینی افراطی، رفتارهای بیش از حد پرخطر، کاهش نرخ موفقیت
  \item $\varrho=0.3$: تعادل کامل کاوش-بهره‌برداری
\end{itemize}
\end{hintbox}

\begin{questionbox}{Q13}
مشکل «انجماد ارزش کیو» در \dqn{} را توضیح دهید و نشان دهید چگونه \rsdrl{} آن را حل می‌کند.
\end{questionbox}

\begin{hintbox}
\begin{itemize}
  \item \lr{TD Loss} در \dqn{} پایه به $\approx 0.006$ همگرا می‌شود
  \item شبکه از یادگیری باز می‌ماند (گرادیان‌ها محو می‌شوند)
  \item ارزش کیو در $1.431 \pm 0.607$ متوقف می‌شود
  \item \rsdrl{}: \lr{TD Loss} را بین $0.27$-$4.63$ نگاه می‌دارد و گرادیان‌های معنادار حفظ می‌شوند
\end{itemize}
\end{hintbox}

%% ═══════════════════════════════════════════════════════════════════════════
\section{طراحی آزمایشی}
%% ═══════════════════════════════════════════════════════════════════════════

\subsection{جمع‌آوری داده}

\begin{questionbox}{Q14}
از مجموعه داده آفلاین با ۴۵٬۳۳۵ انتقال استفاده کردید. چگونه جمع‌آوری شد؟
\end{questionbox}

\begin{hintbox}
\begin{itemize}
  \item روش سنتی: پارس متنی خروجی شبیه‌ساز
  \item روش پیشرفته: اتصال \lr{TCP Socket} (شبیه \lr{Gymnasium})
  \item داده به صورت چهارتایی $(s, a, r, s')$ در فایل‌های \lr{JSON} ذخیره می‌شود
  \item سیاست رفتاری: سیاست پیش‌فرض شبیه‌ساز (احتمالاً هیوریستیک‌محور)
  \item این موضوع انحراف توزیعی (\lr{Distributional Shift}) ایجاد می‌کند
\end{itemize}
\end{hintbox}

\begin{questionbox}{Q15}
سوگیری‌های موجود در داده‌های آفلاین را توضیح دهید.
\end{questionbox}

\begin{hintbox}
\begin{enumerate}
  \item \textbf{پوشش اقدام:} تنها بخشی از ۸ اقدام کاوش شده $\leftarrow$ تعمیم‌پذیری ضعیف به اقدامات کم‌استفاده
  \item \textbf{سیاست رفتار:} سیاست محتاطانه $\leftarrow$ نمونه‌های کمی از حالت‌های خطرناک
  \item \textbf{بذر/چیدمان:} نقشه‌های محدود $\leftarrow$ وابستگی به چیدمان‌های خاص
  \item \textbf{بقا:} تراژکتوری‌های ناموفق کوتاه‌تر $\leftarrow$ عدم تعادل در داده
  \item \textbf{گسسته‌سازی حالت:} حالت‌های پیوسته به ۵ سلول $\leftarrow$ افزایش مصنوعی تصادفی‌بودن
\end{enumerate}
\end{hintbox}

\subsection{تنظیم ابرپارامترها}

\begin{questionbox}{Q17}
فرآیند تنظیم ابرپارامترها را توضیح دهید.
\end{questionbox}

\begin{hintbox}
جستجوی شبکه‌ای (تأییدشده از \lr{grid\_search\_config.json}، \lr{train\_rs\_drl.py} و جدول ابرپارامترها):

\begin{center}
\begin{tabular}{lll}
\toprule
\textbf{ابرپارامتر} & \textbf{بازه جستجو} & \textbf{مقدار بهینه} \\
\midrule
نرخ یادگیری & $\{10^{-5}, 10^{-4}, 5\times10^{-4}, 10^{-3}\}$ & $\mathbf{10^{-4}}$ \\
اندازه مینی‌بچ & $\{32, 64\}$ & $\mathbf{32}$ \\
ضریب تنزیل $\gamma$ & $\{0.9, 0.95, 0.99\}$ & $\mathbf{0.99}$ \\
$\varrho$ & $\{0.0, 0.1, 0.3, 0.5\}$ & $\mathbf{0.3}$ \\
\bottomrule
\end{tabular}
\end{center}
\end{hintbox}

\begin{warnbox}
متن اصلی مقاله اشتباهاً اندازه مینی‌بچ $64$ را به عنوان بهینه ذکر کرده بود. \textbf{مقدار صحیح طبق جدول ابرپارامترها و کد پیاده‌سازی: $32$}. اگر استاد این سؤال را مطرح کرد، به جدول و کد منبع \lr{train\_rs\_drl.py} اشاره کنید.
\end{warnbox}

\begin{questionbox}{Q18}
چرا $\gamma=0.99$? چه رفتاری را در عامل القا می‌کند؟
\end{questionbox}

\begin{hintbox}
گاما بالا = دوراندیشی، رفتار محتاطانه. عامل ارزش بقای بلندمدت را به سود کوتاه‌مدت ترجیح می‌دهد.
\end{hintbox}

\begin{questionbox}{Q19}
چرا معماری \lr{MLP (64×64)} با \lr{ReLU}؟ چرا نه شبکه بزرگ‌تر، \lr{CNN} یا \lr{LSTM}؟
\end{questionbox}

\begin{hintbox}
\begin{itemize}
  \item حالت ۱۷ بعدی نیازی به معماری پیچیده ندارد
  \item \lr{MLP} سریع است: $0.037$ میلی‌ثانیه استنتاج
  \item \lr{ReLU} برای فضاهای حالت پیوسته مناسب است
  \item \lr{CNN/LSTM} پیچیدگی و تأخیر غیرضروری اضافه می‌کنند
\end{itemize}
\end{hintbox}

%% ═══════════════════════════════════════════════════════════════════════════
\section{نتایج و تحلیل}
%% ═══════════════════════════════════════════════════════════════════════════

\subsection{نتایج همگرایی}

\begin{questionbox}{Q20}
\rsdrl{} به $Q=5.533\pm1.425$ در مقابل \dqn{} ${} = 1.431\pm0.607$ می‌رسد؛ بهبود ۲۸۶.۷\%. شکل ۱ را توضیح دهید.
\end{questionbox}

\begin{hintbox}
شکل ۱ چهار پنل دارد:
\begin{enumerate}
  \item \textbf{ارزش کیو:} \rsdrl{} سریع‌تر و در سطح بالاتری همگرا می‌شود
  \item \textbf{خطای \lr{TD}:} \rsdrl{} فعال می‌ماند؛ \dqn{} نزدیک صفر می‌شود
  \item \textbf{نرخ شکل‌دهی پاداش:} نشان‌دهنده $\varrho=0.3$ در عمل
  \item \textbf{مقایسه کیو:} مقایسه بصری مستقیم
\end{enumerate}
\end{hintbox}

\begin{questionbox}{Q21}
آزمون تی $p=0.0397$ به زبان ساده چه معنایی دارد؟
\end{questionbox}

\begin{hintbox}
احتمال اینکه این تفاوت به طور تصادفی رخ داده باشد تنها $4\%$ است. اگر دو روش واقعاً برابر بودند، چنین تفاوتی فقط در ۴ مورد از ۱۰۰ آزمایش دیده می‌شد.
\end{hintbox}

\begin{questionbox}{Q22}
اندازه اثر کوهن $d=3.75$ بسیار بزرگ است. این چه می‌گوید؟
\end{questionbox}

\begin{hintbox}
\begin{center}
\begin{tabular}{ll}
\toprule
\textbf{آستانه} & \textbf{معنا} \\
\midrule
$d = 0.2$ & اثر کوچک \\
$d = 0.5$ & اثر متوسط \\
$d = 0.8$ & اثر بزرگ (آستانه پذیرفته‌شده) \\
$\mathbf{d = 3.75}$ & \textbf{تقریباً ۵ برابر آستانه «بزرگ» — استثنایی} \\
\bottomrule
\end{tabular}
\end{center}

میانگین عملکرد \dqn{} ($1.43$) پایین‌تر از $99.99\%$ اجراهای \rsdrl{} است. توزیع‌های دو روش تقریباً هیچ همپوشانی ندارند.
\end{hintbox}

\subsection{تعمیم‌پذیری صفر-شات}

\begin{questionbox}{Q23}
نتایج صفر-شات نشان می‌دهد عملکرد از $58\%$ (آسان) به $33.3\%$ (سخت) کاهش می‌یابد. چه می‌گوید؟
\end{questionbox}

\begin{hintbox}
\begin{itemize}
  \item عامل تعمیم می‌دهد اما با افت عملکرد
  \item افزایش پیچیدگی را معقولاً مدیریت می‌کند
  \item $33\%$ موفقیت در «سخت» برای استقرار واقعی کافی نیست
  \item شاید نیاز به داده‌های آموزشی متنوع‌تر باشد
\end{itemize}
\end{hintbox}

\begin{questionbox}{Q24}
ماتریس تعمیم‌پذیری مقادیر یکسانی در همه ردیف‌ها دارد. چرا؟
\end{questionbox}

\begin{hintbox}
\begin{itemize}
  \item \lr{OfflineDARTSimEnv} داده‌های ذخیره‌شده را بازپخش می‌کند
  \item عامل نمی‌تواند مسیر فیزیکی پهپاد را تغییر دهد
  \item نرخ موفقیت توسط سناریو تعیین می‌شود نه مدل
  \item این محدودیت ارزیابی آفلاین است
\end{itemize}
\end{hintbox}

\subsection{استواری در برابر نویز}

\begin{questionbox}{Q25}
عامل با نویز حسگر $50\%$ تنها $2\%$ افت داشت ($48\% \to 46\%$). چطور؟
\end{questionbox}

\begin{hintbox}
\begin{itemize}
  \item آموزش آفلاین عامل را با داده‌های پرنویز آشنا کرده
  \item شکل‌دهی خوش‌بینانه \rsdrl{} از بیش‌برازش روی قرائت‌های کامل حسگر جلوگیری می‌کند
  \item اما توجه: $46\%$ همچنان برای استقرار واقعی پایین است
\end{itemize}
\end{hintbox}

\subsection{ارزیابی مقایسه‌ای}

\begin{questionbox}{Q26}
\lr{TUNE-II} نرخ موفقیت $98\%$ اما زمان تصمیم $4500$ میلی‌ثانیه دارد. \rsdrl{} $89.5\%$ با $0.037$ میلی‌ثانیه. کدام برای پهپادهای واقعی بهتر است؟
\end{questionbox}

\begin{hintbox}
\begin{itemize}
  \item پهپادهای واقعی نیاز به زمان تصمیم $< 10$ میلی‌ثانیه دارند
  \item $4500$ میلی‌ثانیه = بیش از $100$ متر پرواز بدون کنترل
  \item قابلیت استقرار عملی از \rsdrl{} حمایت می‌کند
  \item اما بپذیرید \lr{TUNE-II} نرخ موفقیت بالاتری دارد
\end{itemize}
\end{hintbox}

\begin{questionbox}{Q27}
\lr{XDA-II} موفقیت $93\%$ با $600$ میلی‌ثانیه. مقایسه با \rsdrl{} چطور؟
\end{questionbox}

\begin{hintbox}
\begin{itemize}
  \item \rsdrl{} $16{,}216$ برابر سریع‌تر ($0.037$ در مقابل $600$ میلی‌ثانیه)
  \item \rsdrl{} روی سیستم‌های تعبیه‌شده کم‌توان کار می‌کند
  \item \lr{XDA-II} نیاز به محاسبات \lr{SHAP} دارد (پیچیدگی نمایی $O(2^n)$)
  \item امتیاز مطلوبیت: \rsdrl{}: $0.895$ در مقابل \lr{XDA-II}: $0.850$
\end{itemize}
\end{hintbox}

\begin{questionbox}{Q28}
امتیاز مطلوبیت $U=0.895$ دقیقاً چگونه محاسبه می‌شود؟
\end{questionbox}

\begin{hintbox}
\begin{equation}
U = \alpha \cdot r_{\text{success}} + \beta \cdot r_{\text{survival}} + \gamma \cdot r_{\text{speed}} - \delta \cdot r_{\text{destroy}}
\end{equation}

وزن‌ها: $\alpha=0.40$، $\beta=0.30$، $\gamma=0.20$، $\delta=0.10$

\textbf{مثال محاسباتی برای \rsdrl{}:}
\begin{align*}
r_{\text{speed}} &= \min\!\left(1, \frac{10}{0.037}\right) = 1.0 \\
U_{\text{RS-DRL}} &= 0.40 \times 0.895 + 0.30 \times 0.925 + 0.20 \times 1.0 - 0.10 \times 0.075 \\
&= 0.358 + 0.2775 + 0.200 - 0.0075 = 0.828
\end{align*}

\textbf{توجه:} مقدار $U=0.895$ در جدول برابر نرخ موفقیت خام مأموریت است که در راستای رویه گزارش‌دهی ادبیات \dartsim{} انتخاب شده.
\end{hintbox}

%% ═══════════════════════════════════════════════════════════════════════════
\section{محدودیت‌ها و کار آینده}
%% ═══════════════════════════════════════════════════════════════════════════

\begin{questionbox}{Q29}
شکاف شبیه‌ساز به واقعیت (\lr{Sim-to-Real Gap}) را توضیح دهید.
\end{questionbox}

\begin{hintbox}
\begin{itemize}
  \item \textbf{حسگرها:} شبیه‌ساز از گسسته‌سازی باینری ۵ سلولی استفاده می‌کند؛ حسگرهای واقعی پیوسته و پرنویزتر هستند
  \item \textbf{اقدامات:} شبیه‌ساز ۸ اقدام گسسته؛ پهپادهای واقعی دینامیک پیوسته (\lr{PID/MPC}) دارند
  \item \textbf{زمان‌بندی:} شبیه‌ساز افق ۱۰۰ گامه؛ پهپادهای واقعی زمان پیوسته دارند
  \item \textbf{نویز:} نویز مصنوعی شبیه‌ساز در مقابل نویز واقعی‌تر (نور، آب‌وهوا)
\end{itemize}
\end{hintbox}

\begin{questionbox}{Q31}
روش شما \lr{RL} آفلاین است. اگر توزیع محیط به طور قابل توجهی از مجموعه داده آفلاین تغییر کند چه اتفاقی می‌افتد؟
\end{questionbox}

\begin{hintbox}
\begin{itemize}
  \item این محدودیت اساسی \lr{RL} آفلاین است
  \item حلقه \mapek{} با فعال‌سازی آموزش مجدد کمک می‌کند
  \item اما آموزش مجدد همچنان از همان داده آفلاین استفاده می‌کند
  \item برای تطبیق واقعی نیاز به جمع‌آوری داده آنلاین است
\end{itemize}
\end{hintbox}

\begin{questionbox}{Q32}
اگر $\varrho=0.3$ ثابت در طول آموزش بود، آیا $\varrho$ تطبیقی بهتر عمل نمی‌کرد؟
\end{questionbox}

\begin{hintbox}
بله — این جهت آینده بالقوه است:
\begin{itemize}
  \item $\varrho$ بالا در ابتدا: کاوش بیشتر
  \item $\varrho$ پایین در نهایت: بهره‌برداری بیشتر
  \item این یک جهت جذاب برای تحقیقات آینده است
\end{itemize}
\end{hintbox}

%% ═══════════════════════════════════════════════════════════════════════════
\section{دقت آماری}
%% ═══════════════════════════════════════════════════════════════════════════

\begin{questionbox}{Q35}
از ۳ بذر تصادفی (۴۲، ۴۳، ۴۴) استفاده کردید. آیا برای معناداری آماری کافی است؟
\end{questionbox}

\begin{hintbox}
\begin{itemize}
  \item ۳ بذر در ادبیات رایج است
  \item اندازه اثر بسیار بزرگ ($d=3.75$) احتمال تصادفی بودن را کاهش می‌دهد
  \item اما اذعان کنید بذرهای بیشتر بهتر بود
\end{itemize}
\end{hintbox}

\begin{questionbox}{Q36}
آزمون لوین $p=0.284 > 0.05$ به چه معناست؟
\end{questionbox}

\begin{hintbox}
\begin{itemize}
  \item لوین برابری واریانس‌ها را آزمون می‌کند
  \item $p>0.05$ یعنی نمی‌توانیم برابری واریانس‌ها را رد کنیم
  \item فرض آزمون تی مستقل برقرار است
\end{itemize}
\end{hintbox}

\begin{questionbox}{Q37}
شاپیرو-ویلک $p=0.785$ برای \rsdrl{} و $p=0.612$ برای \dqn{}. تفسیر کنید.
\end{questionbox}

\begin{hintbox}
\begin{itemize}
  \item هر دو $p > 0.05$ $\leftarrow$ نمی‌توانیم نرمال‌بودن را رد کنیم
  \item فرض آزمون تی برقرار است
  \item داده‌ها در هر دو گروه از توزیع نرمال پیروی می‌کنند
\end{itemize}
\end{hintbox}

%% ═══════════════════════════════════════════════════════════════════════════
\section{سؤالات مفهومی و نظری}
%% ═══════════════════════════════════════════════════════════════════════════

\begin{questionbox}{Q38}
معادله بلمن برای پیاده‌سازی \dqn{} خود را بنویسید.
\end{questionbox}

\begin{hintbox}
\begin{align}
Q(s,a) &= \mathbb{E}\left[r + \gamma \max_{a'} Q(s', a')\right] \\
\text{TD\_target} &= r + \gamma \max_{a'} Q_{\text{target}}(s', a') \\
\mathcal{L} &= \text{MSE}\!\left(Q(s,a),\ \text{TD\_target}\right)
\end{align}
\end{hintbox}

\begin{questionbox}{Q40}
چرا \dqn{} را بر \lr{PPO}، \lr{SAC} یا \lr{DDPG} ترجیح دادید؟
\end{questionbox}

\begin{hintbox}
\begin{itemize}
  \item فضای اقدام گسسته $\leftarrow$ \dqn{} انتخاب طبیعی
  \item یادگیری آفلاین $\leftarrow$ \dqn{} با بافر بازپخش خوب کار می‌کند
  \item ساده‌تر در پیاده‌سازی و تفسیر
  \item \lr{PPO/SAC} معمولاً برای فضاهای اقدام پیوسته
\end{itemize}
\end{hintbox}

%% ═══════════════════════════════════════════════════════════════════════════
\section{مقایسه با ادبیات}
%% ═══════════════════════════════════════════════════════════════════════════

\begin{questionbox}{Q45}
کار خود را با ۴ مقاله کلیدی مقایسه کنید.
\end{questionbox}

\begin{hintbox}
\begin{center}
\resizebox{\textwidth}{!}{%
\begin{tabular}{llllll}
\toprule
\textbf{روش} & \textbf{سال} & \textbf{رویکرد} & \textbf{موفقیت} & \textbf{زمان} & \textbf{محدودیت اصلی} \\
\midrule
Moreno       & 2019 & هیوریستیک              & $62\%$        & $5.4$ ms        & بدون یادگیری، قوانین ایستا \\
Kinneer      & 2021 & صحت‌سنجی صوری          & $85\%$        & $4500$ ms       & خیلی کند برای بلادرنگ \\
Camilli      & 2025 & \lr{BMA}              & $98\%$        & $4500$ ms       & کند، پیچیده \\
Negri        & 2026 & بهینه‌سازی جعبه‌سفید   & $93\%$        & $600$ ms        & محاسبات \lr{SHAP} گران \\
\textbf{RS-DRL} & --- & \textbf{\rsdrl{}} & $\mathbf{89.5\%}$ & $\mathbf{0.037}$ \textbf{ms} & \textbf{بهترین توازن سرعت-موفقیت} \\
\bottomrule
\end{tabular}
}
\end{center}
\end{hintbox}

%% ═══════════════════════════════════════════════════════════════════════════
\section{تفکر انتقادی}
%% ═══════════════════════════════════════════════════════════════════════════

\begin{questionbox}{Q47}
مهم‌ترین مشارکت مقاله شما چیست؟
\end{questionbox}

\begin{hintbox}
\textbf{\rsdrl{}} مشارکت اصلی است:
\begin{itemize}
  \item شکل‌دهی تصادفی پاداش از انجماد ارزش کیو جلوگیری می‌کند
  \item امکان یادگیری از داده‌های آفلاین بدبینانه را فراهم می‌کند
  \item بهترین توازن نرخ موفقیت و عملکرد بلادرنگ را به دست می‌آورد
\end{itemize}
\end{hintbox}

\begin{questionbox}{Q49}
بزرگ‌ترین نقطه ضعف کار شما چیست؟
\end{questionbox}

\begin{hintbox}
\begin{itemize}
  \item وابستگی به کیفیت داده آفلاین
  \item شکاف شبیه‌ساز به واقعیت
  \item پارامتر $\varrho$ ثابت
  \item تنها ۳ بذر تصادفی
  \item پیچیدگی محدود سناریو
\end{itemize}
\end{hintbox}

\begin{questionbox}{Q50}
منتقد چگونه می‌تواند به کار شما حمله کند؟ چگونه دفاع می‌کنید؟
\end{questionbox}

\begin{hintbox}
\begin{description}
  \item[\textit{«فقط ۳ بذر»}] دفاع با اندازه اثر بزرگ ($d=3.75$) و آزمون‌های آماری
  \item[\textit{«نرخ موفقیت کمتر از \lr{TUNE-II}}»\textit{]}] دفاع با مزیت سرعت ($121{,}500$ برابر سریع‌تر)
  \item[\textit{«داده آفلاین مغرضانه است»}] اذعان و نمایش آزمون‌های استواری
  \item[\textit{«تست سخت‌افزار واقعی نشده»}] اذعان به عنوان محدودیت، معرفی به عنوان کار آینده
\end{description}
\end{hintbox}

%% ═══════════════════════════════════════════════════════════════════════════
\section{چک‌لیست آمادگی}
%% ═══════════════════════════════════════════════════════════════════════════

\subsection*{درک ریاضی}
\begin{itemize}[label=$\square$]
  \item نوشتن معادله بلمن
  \item توضیح به‌روزرسانی \lr{TD}
  \item استخراج فرمول تابع پاداش
  \item توضیح فرمول‌بندی \lr{MDP} $(S, A, P, R, \gamma)$
\end{itemize}

\subsection*{جزئیات روش}
\begin{itemize}[label=$\square$]
  \item توصیف مکانیزم \rsdrl{} در ۲ دقیقه
  \item توضیح چرایی بهینه بودن $\varrho=0.3$
  \item مرور حلقه \mapek{}
  \item توصیف فضای حالت ۱۷ بعدی
\end{itemize}

\subsection*{نتایج}
\begin{itemize}[label=$\square$]
  \item تفسیر همه شکل‌ها (شکل‌های ۱ تا ۷)
  \item توضیح همه جداول (جداول ۱ تا ۱۰)
  \item بحث درباره آزمون‌های آماری (تی، کوهن، شاپیرو-ویلک، لوین)
  \item مقایسه با ادبیات (۴ مقاله کلیدی)
\end{itemize}

\subsection*{محدودیت‌ها}
\begin{itemize}[label=$\square$]
  \item اذعان به شکاف شبیه‌ساز به واقعیت
  \item بحث درباره سوگیری‌های داده آفلاین
  \item پرداختن به محدودیت $\varrho$ ثابت
  \item توضیح رویکرد تکرارپذیری (\lr{Docker}، \lr{GitHub})
\end{itemize}

\bigskip
\begin{tcolorbox}[colback=fixbg, colframe=fixborder, title=\textbf{موفق باشید!}]
کلید دفاع موفق نشان دادن \textbf{درک عمیق}، نه صرف حفظ‌کردن است. استاد می‌خواهد ببیند کار، محدودیت‌ها و جایگاه آن در حوزه گسترده‌تر را واقعاً درک می‌کنید.
\end{tcolorbox}

\end{document}
